\chapter{总结}

辐射传输方程以相空间（包含空间位置、速度方向及能量）为自变量，能够综合描述粒子的散射、吸收和源项等物理过程，从而在统一的框架下精确刻画剂量沉积分布。然而，直接求解这类高维输运方程面临巨大的数值挑战：一方面需要处理复杂的边界和源条件（如多方向外照射束、内部放射源等），另一方面还需满足临床应用对计算效率和稳定性的严格要求。尽管蒙特卡洛方法在理论上能够逼近真实的输运过程，并被视为剂量计算的“金标准”，但其高昂的计算成本限制了其在常规临床流程中的广泛应用。

因此，如何在降低计算代价的同时，既保留输运方程的物理精度，又兼顾临床应用的实时性需求，已成为当前放疗剂量计算领域亟待解决的关键问题。

针对上述挑战，本文提出了一种用于求解辐射传输方程的神经算子框架——DeepRTE。该框架将注意力机制与算子学习策略有机结合，能够直接建立从入射边界条件、散射系数及散射核等输入参数到方程解的映射关系。其主要创新点如下：
\begin{enumerate}
    \item \textbf{参数高效性}：通过在网络结构中显式编码物理规律，DeepRTE 在参数量远少于纯数据驱动大模型的前提下，依然可以达到甚至超过后者的精度；
    \item \textbf{零样本泛化能力}：在保持网络参数不变的情况下，DeepRTE 能够对训练阶段未曾见过的边界条件进行稳健预测，展现出良好的零样本外推能力；
    \item \textbf{可解释性}：借助物理引导的网络设计，模型整体保持线性、物理含义清晰的算子形式，便于从物理视角理解网络行为。
\end{enumerate}

数值实验表明，DeepRTE 具有显著的迁移学习与零样本能力。尽管模型只在特定形式的入射边界条件上进行了训练，当将其直接应用到全新的边界条件时，依然能够取得较高精度：预测解与数值真解高度一致，各类误差指标保持在较低水平。这表明 DeepRTE 并非简单记忆训练数据，而是有效捕捉并内化了辐射传输方程所蕴含的基本物理规律。

更进一步地，DeepRTE 在完全未见过的边界条件下仍能给出准确预测，其零样本表现尤为突出。在不进行任何额外训练或微调的前提下，模型即可对全新场景下的辐射传输问题给出较为可靠的数值解。这一现象从侧面印证了物理先验与深度学习框架的有机结合，为构建具有物理可信度的神经算子提供了可行路径。

总体而言，DeepRTE 展示了将深度学习方法与物理定律相融合以求解复杂科学与工程问题的潜力。一方面，模型能够在无需反复求解传统数值离散系统的情况下，快速给出辐射传输方程的高精度近似解，大幅提升了计算效率；另一方面，得益于物理引导的网络结构设计，模型在面对未见过的条件和分布外场景时仍具有良好的稳健性与泛化能力。未来，可以在更高维度、更复杂几何域以及更一般的散射模型下进一步拓展 DeepRTE 的应用范围，并结合自适应采样、高阶数值积分等技术，进一步提升整体精度与效率，为科学计算领域中其他复杂偏微分方程与积分方程问题提供一种具有参考价值的通用范式。